% GENERATED FILE. Edit canonical lessons, then run npm run generate:books.
% canonical-source-hash: fnv1a64:946941afd499920e
% canonical-lessons: PT-C369-condicao, PT-C369-desde-que, PT-C369-contanto-que, PT-C369-caso, PT-C369-consoante

\chapter{Com uma condição}
\label{ch:pt-com-uma-condicao}
\hlchaptermodality{369}

\begin{chapteropening}
\textbf{By the end of this chapter:} \emph{I can set conditions: provided that, as long as, in case, and depending on something.}

Complete the last lesson of chapter 369: I can set conditions: provided that, as long as, in case, and depending on something.
\end{chapteropening}

\section[a condição]{a condição — the condition}

\label{lesson:PT-C369-condicao}



\begin{quote}
Before the new one: say the Portuguese for to presume, then the Portuguese for plausible.
\end{quote}

\subsection*{You'll want to know: a condição}
\textbf{a condição} — \textquotedblleft{}the condition\textquotedblright{}.

\begin{quote}
Vou contigo, com uma condição: voltamos cedo.
\end{quote}

\begin{quote}
\emph{I will go with you on one condition: we come back early.}
\end{quote}

\emph{Com a condição de} + infinitive. \emph{Condições de trabalho} — working conditions.

\begin{cousinweb}
\emph{condição} is a learned borrowing from Latin \emph{condiciō} \textquotedblleft{}an agreement, a condition\textquotedblright{} (EtymDB, Lewis \& Short), often spelt \emph{conditiō} in later Latin, built on \emph{condīcere} \textquotedblleft{}to talk together, to agree on\textquotedblright{}, from \emph{dīcere} \textquotedblleft{}to say\textquotedblright{}.
\end{cousinweb}

\begin{grammarlens}[title={What you've built}]
One more word for telling, explaining and describing at B1.
\end{grammarlens}

\subsection*{Your turn}
\begin{itemize}
\raggedright
  \item \emph{Say it:} \emph{a condição}
  \item \emph{Say it:} \emph{a condição}, once more
  \item \emph{Say it:} a sentence of your own with \emph{a condição}
  \item \emph{Recall:} say the Portuguese for to presume, then the Portuguese for plausible, then say \emph{a condição} again
\end{itemize}

\begin{tcolorbox}[breakable,colback=teal!4,colframe=teal!35!black,title={Before you move on}]
What does a condição mean? (\textquotedblleft{}The condition\textquotedblright{}.) Say it once more.
\end{tcolorbox}
\section[desde que]{desde que — provided that, as long as; since (in time)}

\label{lesson:PT-C369-desde-que}



\begin{quote}
Before the new one: say the Portuguese for plausible, then the Portuguese for the condition.
\end{quote}

\subsection*{You'll want to know: desde que}
\textbf{desde que} — \textquotedblleft{}provided that, as long as; since (in time)\textquotedblright{}.

\begin{quote}
Podes ir à festa desde que voltes antes da meia-noite.
\end{quote}

\begin{quote}
\emph{You can go to the party as long as you come back before midnight.}
\end{quote}

With the subjunctive it sets a condition: \emph{desde que voltes}. With the indicative it means since: \emph{desde que cheguei}.

\begin{cousinweb}
\emph{desde que} is built from \emph{desde} \textquotedblleft{}since\textquotedblright{}, usually traced to Latin \emph{dē} + \emph{ex} + \emph{dē}, all meaning \textquotedblleft{}from\textquotedblright{} (Lewis \& Short), and \emph{que}.
\end{cousinweb}

\begin{grammarlens}[title={What you've built}]
One more word for telling, explaining and describing at B1.
\end{grammarlens}

\subsection*{Your turn}
\begin{itemize}
\raggedright
  \item \emph{Say it:} \emph{desde que}
  \item \emph{Say it:} \emph{desde que}, once more
  \item \emph{Say it:} a sentence of your own with \emph{desde que}
  \item \emph{Recall:} say the Portuguese for plausible, then the Portuguese for the condition, then say \emph{desde que} again
\end{itemize}

\begin{tcolorbox}[breakable,colback=teal!4,colframe=teal!35!black,title={Before you move on}]
What does desde que mean? (\textquotedblleft{}Provided that, as long as; since (in time)\textquotedblright{}.) Say it once more.
\end{tcolorbox}
\section[contanto que]{contanto que — as long as, on condition that}

\label{lesson:PT-C369-contanto-que}



\begin{quote}
Before the new one: say the Portuguese for the condition, then the Portuguese for provided that.
\end{quote}

\subsection*{You'll want to know: contanto que}
\textbf{contanto que} — \textquotedblleft{}as long as, on condition that\textquotedblright{}.

\begin{quote}
Empresto-te o carro, contanto que tenhas cuidado.
\end{quote}

\begin{quote}
\emph{I will lend you the car, as long as you are careful.}
\end{quote}

Takes the subjunctive; a little more formal than \emph{desde que}.

\begin{cousinweb}
\emph{contanto que} is built from \emph{com} \textquotedblleft{}with\textquotedblright{}, \emph{tanto} \textquotedblleft{}so much\textquotedblright{}, inherited from Latin \emph{tantum} \textquotedblleft{}so much\textquotedblright{} (Lewis \& Short), and \emph{que}: \textquotedblleft{}with so much as\textquotedblright{}.
\end{cousinweb}

\begin{grammarlens}[title={What you've built}]
One more word for telling, explaining and describing at B1.
\end{grammarlens}

\subsection*{Your turn}
\begin{itemize}
\raggedright
  \item \emph{Say it:} \emph{contanto que}
  \item \emph{Say it:} \emph{contanto que}, once more
  \item \emph{Say it:} a sentence of your own with \emph{contanto que}
  \item \emph{Recall:} say the Portuguese for the condition, then the Portuguese for provided that, then say \emph{contanto que} again
\end{itemize}

\begin{tcolorbox}[breakable,colback=teal!4,colframe=teal!35!black,title={Before you move on}]
What does contanto que mean? (\textquotedblleft{}As long as, on condition that\textquotedblright{}.) Say it once more.
\end{tcolorbox}
\section[caso]{caso — in case, if (+ subjunctive); the case}

\label{lesson:PT-C369-caso}



\begin{quote}
Before the new one: say the Portuguese for provided that, then the Portuguese for as long as.
\end{quote}

\subsection*{You'll want to know: caso}
\textbf{caso} — \textquotedblleft{}in case, if (+ subjunctive); the case\textquotedblright{}.

\begin{quote}
Caso chova, o jantar é dentro de casa.
\end{quote}

\begin{quote}
\emph{If it rains, dinner is indoors.}
\end{quote}

\emph{Caso} + present subjunctive: \emph{caso chova}. As a noun, \emph{o caso} is the case: \emph{neste caso}.

\begin{cousinweb}
\emph{caso} comes from Latin \emph{cāsus} \textquotedblleft{}a falling, an event, a chance\textquotedblright{} (EtymDB, Lewis \& Short), built on \emph{cadere} \textquotedblleft{}to fall\textquotedblright{}: what befalls you.
\end{cousinweb}

\begin{grammarlens}[title={What you've built}]
One more word for telling, explaining and describing at B1.
\end{grammarlens}

\subsection*{Your turn}
\begin{itemize}
\raggedright
  \item \emph{Say it:} \emph{caso}
  \item \emph{Say it:} \emph{caso}, once more
  \item \emph{Say it:} a sentence of your own with \emph{caso}
  \item \emph{Recall:} say the Portuguese for provided that, then the Portuguese for as long as, then say \emph{caso} again
\end{itemize}

\begin{tcolorbox}[breakable,colback=teal!4,colframe=teal!35!black,title={Before you move on}]
What does caso mean? (\textquotedblleft{}In case, if (+ subjunctive); the case\textquotedblright{}.) Say it once more.
\end{tcolorbox}
\section[consoante]{consoante — depending on, according to; the consonant}

\label{lesson:PT-C369-consoante}



\begin{quote}
Before the new one: say the Portuguese for as long as, then the Portuguese for in case.
\end{quote}

\subsection*{You'll want to know: consoante}
\textbf{consoante} — \textquotedblleft{}depending on, according to; the consonant\textquotedblright{}.

\begin{quote}
Vamos à praia ou ao cinema, consoante o tempo.
\end{quote}

\begin{quote}
\emph{We will go to the beach or to the cinema, depending on the weather.}
\end{quote}

\emph{Consoante} + noun: \emph{consoante o preço}. As a noun, \emph{a consoante} is a consonant.

\begin{cousinweb}
\emph{consoante} is inherited from Old Portuguese \emph{consõante} (EtymDB), from Latin \emph{cōnsonāns} \textquotedblleft{}sounding together, agreeing\textquotedblright{} (Lewis \& Short), on \emph{sonāre} \textquotedblleft{}to sound\textquotedblright{}: what agrees with something goes according to it.
\end{cousinweb}

\begin{grammarlens}[title={What you've built}]
One more word for telling, explaining and describing at B1.
\end{grammarlens}

\subsection*{Your turn}
\begin{itemize}
\raggedright
  \item \emph{Say it:} \emph{consoante}
  \item \emph{Say it:} \emph{consoante}, once more
  \item \emph{Say it:} a sentence of your own with \emph{consoante}
  \item \emph{Recall:} say the Portuguese for as long as, then the Portuguese for in case, then say \emph{consoante} again
\end{itemize}

\begin{tcolorbox}[breakable,colback=teal!4,colframe=teal!35!black,title={Before you move on}]
What does consoante mean? (\textquotedblleft{}Depending on, according to; the consonant\textquotedblright{}.) Say it once more.
\end{tcolorbox}
